%------------------------------------------------------------------------Begin preamble------------------------------------------------------------------------
\documentclass[12pt]{article}
% Uncomment these two lines if you want to use Times New Roman (needs XeLaTeX)
% \usepackage{fontspec}
% \setmainfont{Times New Roman}

\usepackage{amsmath} % For math
\usepackage{amssymb} % For more math
\usepackage{fancyhdr} % For fancy headers and footers
\usepackage{fancyvrb} % For writing blocks of code verbatim (like LaTeX code)
\usepackage{geometry} % For manipulating margins and meta doc stuff
\usepackage{graphicx} % Required for inserting images
\usepackage{hyperref} % For linking stuff
\usepackage{cleveref} % Better cross refencing % Must be loaded after hyperref
\usepackage{lastpage} % For doing "Page n of |n|"
\usepackage{minted} % For color coding code sections
\usepackage{pgfplots} % For adding plots and graphs
\usepackage[indent=00pt]{parskip} % Get rid of beginning indents on paragraphs and set spacing between paragraphs
\usepackage{siunitx} % For adding siunits inside math zones
\usepackage[most]{tcolorbox} % For inserting color boxes
\usepackage{tikz}
\usepackage{titling} % For vertically centering title, author, and date
\usepackage{tocloft} % For cooler toc
\usepackage{xcolor} % For coloring text
\usepackage{witharrows} % For cool arrows for pointing at stuff

%%%%%%%%% Begin tocloft setup %%%%%%%%%%%%%%
\renewcommand\cftsecfont{\normalfont}
\renewcommand\cftsecpagefont{\normalfont}
\renewcommand{\cftsecleader}{\cftdotfill{\cftsecdotsep}}
\renewcommand\cftsecdotsep{\cftdot}
\renewcommand\cftsubsecdotsep{\cftdot}
%%%%%%%%% End tocloft %%%%%%%%%%%%%%

%%%%%%%%%% Begin hyperref setup %%%%%%%%%%%%%%%%%%%
\hypersetup{
    colorlinks=true,
    linkcolor=blue,
    urlcolor=blue,
}
%%%%%%%%%% End hyperref setup %%%%%%%%%%%%%%%%%%%

%%%%%%% Begin pgfplotsset setup %%%%%%%
\pgfplotsset{compat=1.18} % Compiler be bugging
%%%%%%% End pgfplotsset setup %%%%%%%

%%%%%%% Begin fancyhdr setup %%%%%%%
\pagestyle{fancy}
\renewcommand{\footrulewidth}{0.4pt} % default is 0pt
\lhead{Linear Algebra} % Top left header
\chead{} % Top center header
\rhead{} % Top right header

\lfoot{Brandon J. T. Noguera} % Bottom left footer
\cfoot{} % Bottom center footer
\rfoot{Page \thepage\ of \pageref*{LastPage}} % Bottom right footer
% doing \pageref*{some_url} turns off the hyperlink capability for some_url
%%%%%%% End fancyhdr setup %%%%%%%

%%%%%%% Begin titlingpage setup %%%%%%%
\renewcommand\maketitlehooka{\null\mbox{}\vfill}
\renewcommand\maketitlehookd{\vfill\null}
%%%%%% End titlingpage setup %%%%%%%
%------------------------------------------------------------------------End preamble------------------------------------------------------------------------

\title{\textbf{Linear Algebra}}
\author{Brandon Jose Tenorio Noguera}
\date{}
\begin{document}

\begin{titlingpage}
\maketitle
\end{titlingpage}

\tableofcontents
\newpage

\section{Sources}\label{srcs}

\begin{tcolorbox}[title=Sources]
These notes were taken from
\href{https://ocw.mit.edu/courses/res-18-010-a-2020-vision-of-linear-algebra-spring-2020/resources/a-new-way-to-start-linear-algebra/}{MIT's OCW: A vision of linear algebra} and \href{https://ocw.mit.edu/courses/18-06sc-linear-algebra-fall-2011/}{MIT's OCW: Linear Algebra}
\end{tcolorbox}

\section{The column space of a matrix}\label{p1}
\begin{tcolorbox}[title=Sources]
\href{https://ocw.mit.edu/courses/res-18-010-a-2020-vision-of-linear-algebra-spring-2020/resources/the-column-space-of-a-matrix/}{The Column Space of a Matrix by Professor Gilbert Strang}
\end{tcolorbox}

\subsection{$A = \mathbf{CR}$}\label{}
\begin{align*}\label{}
  A_0 &= 
 \begin{bmatrix}
   1 & 3 & 2 \\
   4 & 12 & 8 \\
   2 & 6 & 4
 \end{bmatrix} 
 \\
  r_1 &= \begin{bmatrix}
    1 & 3 & 2
  \end{bmatrix}
  \\
    r_2 &= \begin{bmatrix}
    4 & 12 & 8
  \end{bmatrix}
  \\
      r_3 &= \begin{bmatrix}
    2 & 6 & 4
  \end{bmatrix}
  \\
  c_1 &= \begin{bmatrix}
    1 \\
    4\\
    2
  \end{bmatrix}
  \\
    c_2 &= \begin{bmatrix}
    3\\
    12\\
    6
    \end{bmatrix}
    \\
      c_3 &= \begin{bmatrix}
      2\\
      8\\
      4
      \end{bmatrix}
\end{align*}

Notice that $r_2 = 4(r_1) \text{ and } r_3 = 2(r_2)$

Notice also that $c_2 = 3 (c_1) \text{ and } c_3 = 2 (c_1) $

\begin{align*}
  A_1 &= \begin{bmatrix}
    1 & 4 & 2\\
    4 & 1 & 3\\
    5 & 5 & 5
  \end{bmatrix}
\end{align*}

Notice that $ r_1 + r_2 = r_3 $

These are symmetrical matrices, that's why they start with $ S $
\begin{equation*}
    S_2 = \begin{bmatrix}
      2 & -1\\
      -1 & 2
    \end{bmatrix}
    \ 
    S_3 = \begin{bmatrix}
      1 & -1 & 0\\
      -1 & 2 & -1\\
      0 & -1 & 1
    \end{bmatrix}
    \ 
    S_4 = \begin{bmatrix}
      2 & -1 & 0\\
      -1 & 2 & -1\\
      0 & -1 & 2
    \end{bmatrix}
\end{equation*}

These are \emph{orthogonal} matrices (this just means perpendicular)

Orthogonal matrices tend to express a rotation (they rotate the plane)

\begin{equation*}
  Q_5 = \begin{bmatrix}
    \cos^{}\left(\theta\right) & -\sin^{}\left(\theta\right)\\
    \sin^{}\left(\theta\right) & \cos^{}\left(\theta\right)
  \end{bmatrix}
\end{equation*}

\section{The column space of $A $}\label{}

\begin{align*}
  A \vec{x} &= \begin{bmatrix}
    1 & 4 & 5\\
    3 & 2 & 5\\
    2 & 1 & 3
  \end{bmatrix} \cdot \begin{bmatrix}
  x_1\\
  x_2\\
  x_3
  \end{bmatrix}
  =
  \begin{bmatrix}
  1\\
  3\\
  2
  \end{bmatrix} \cdot x_1
  + 
  \begin{bmatrix}
  4\\
  2\\
  1
  \end{bmatrix} \cdot x_2
  +
  \begin{bmatrix}
    5\\
    5\\
    3
  \end{bmatrix} \cdot x_3
  \\
  &= 
  \text{one linear combination of the columns of }A
\end{align*}

"Linear" because nothing is being squared or cubed or something\\
"Combination" because we're putting them together

\subsection{All combinations of A}\label{}

The column space of $A = \text{all vector } A \mathbf{x} $ 
\end{document}
